\section{Spin systems}
\label{sec:setup}

This section fixes the spin-system setup and notation used throughout the paper.
We describe the dynamics through a neighbourhood-dependent update rate and
refresh parameter, which directly extends the usual KCSM formulation.

Let~$\Lambda$ be a countable site set. For each~$i\in\Lambda$, fix a finite
neighbourhood~$N(i)\subseteq\Lambda\setminus\cb{i}$ and assume
\begin{equation*}
\sup_{i\in\Lambda}\abs{N(i)}<\infty.
\end{equation*}
The neighbourhood relation may be directed. \gm{not clear what is meant with directed neighborhood relation. maybe just leave this comment out.} Write~$N_*(i)=N(i)\cup\cb{i}$
and, for~$A\subseteq\Lambda$,
\begin{equation*}
N_*(A)=\bigcup_{i\in A}N_*(i),
\qquad
B(A,0)=A,
\qquad
B(A,n+1)=N_*\bb{B(A,n)}.
\end{equation*}
Thus~$B(A,n)$ is the directed ball reached from~$A$ in at most~$n$ steps. For
$i\in\Lambda$, write~$B(i,n)=B(\cb{i},n)$. We say that the neighbourhood graph
has polynomial growth if, for some~$K,D<\infty$,
\begin{equation}
\label{eq:polynomial-growth}
\abs{B(i,n)}\leq K(1+n)^D
\qquad
\text{for every }i\in\Lambda\text{ and }n\in\N.
\end{equation}


\begin{definition}[Orientability]
\label{def:orientability}
A set~$V\subseteq\Lambda$ is \emph{autonomous} if~$N_*(V)=V$.
The neighbourhood graph is \emph{orientable} if every finite set~$A$ is
contained in a finite set~$\bar A$ for which there is
$U\subseteq\Lambda\setminus\bar A$ such that
\begin{equation*}
N_*(U)=U,
\qquad
N_*(\bar A\cup U)=\bar A\cup U.
\end{equation*}
\end{definition}
For example, if~$\Lambda=\Z$ and~$N(i)\subseteq\cb{j:j>i}$, then for
every nonempty finite~$A$ we may take
\begin{equation*}
\bar A=[\min A,\max A]\cap\Z,
\qquad U=\cb{j\in\Z:j>\max A}.
\end{equation*}
Both~$U$ and~$\bar A\cup U$ are autonomous, so the graph is orientable.
\par

Put~$\Omega=\cb{0,1}^{\Lambda}$. A local function depends on finitely many
coordinates. For~$x\in\cb{0,1}$, let~$\eta^{i,x}$ be the configuration obtained
from~$\eta$ \gm{$\eta$ not defined need $\eta \in \Omega$} by setting the spin at~$i$ equal to~$x$, and let~$\eta^i$ be obtained
by flipping it. For~$p\in[0,1]$, the Bernoulli refresh operator at~$i$ is
\begin{align*}
E_{i,p}f(\eta)
&=
pf(\eta^{i,1})+(1-p)f(\eta^{i,0}),
\\
\bb{E_{i,p}-I}f(\eta)
&=
\bb{p\bb{1-\eta(i)}+(1-p)\eta(i)}
\bb{f(\eta^i)-f(\eta)}.
\end{align*}
Thus a refresh sets the spin to~$1$ with probability~$p$ and to~$0$ with
probability~$1-p$.

\gm{Make clear that in contrast to the classical approach to spin systems (see Ligget) we use a special form of the operator that is particulary adapted to, and useful for, studying KCSMs.}

A spin system is specified by functions \gm{what is ~$i$? I guess~$i\int \Lambda$.  }
\begin{equation*}
\lambda_i:\Omega\to[0,\infty),
\qquad
\rho_i:\Omega\to[0,1],
\end{equation*}
each depending only on the spins in~$N(i)$, with
$\sup_{i,\eta}\lambda_i(\eta)<\infty$. Its generator is \gm{we need to specify what is ~$f$ local functions...?}
\begin{equation}
\label{eq:spin-generator}
\Lcal f(\eta)
=
\sum_{i\in\Lambda}
\lambda_i(\eta)\bb{E_{i,\rho_i(\eta)}-I}f(\eta).
\end{equation}
\gm{Obvious question we need to address (maybe in a remark): this generator differs from standard definition of spin systems in the sense of Ligget for example. are the definitions equivalent?} At rate~$\lambda_i(\eta)$, the spin at~$i$ is refreshed to~$1$ with
probability~$\rho_i(\eta)$ and to~$0$ otherwise. The value of~$\rho_i(\eta)$ is
irrelevant when~$\lambda_i(\eta)=0$.  \gm{maybe better to do an example environment and give the examples of IPS we use in this manuscript explicitly.} For a standard KCSM with Bernoulli refresh
parameter~$p$, one has~$\rho_i\equiv p$ and~$\lambda_i$ is the kinetic
constraint. More generally, here both the update rate and the refresh parameter
may depend on the neighbourhood.

\gm{Next paragraph is misleading it sounds like every spin system can be written in our paramaterization, which is wrong.} This parametrization is equivalent to the usual flip-rate description. If
$c_i$ denotes the flip rate in the conventional representation
$\Lcal f(\eta)=\sum_i c_i(\eta)\bb{f(\eta^i)-f(\eta)}$, then \gm{second and third formula need explanation and/or reminder that $\lambda_i$ and $p_i$ do not depend on the spin at site~$i$ therefore~$\lambda_i(\eta) = \lambda_i(\eta^{i,0}) = \lambda_i(\eta^{i,1})$ and the same for $p_i$, so second formula follows from first formula. }
\begin{align*}
c_i(\eta)
&=
\lambda_i(\eta)
\bb{
\bb{1-\eta(i)}\rho_i(\eta)
+
\eta(i)\bb{1-\rho_i(\eta)}
},
\\
\lambda_i(\eta)
&=
c_i(\eta^{i,0})+c_i(\eta^{i,1}),
\\
\rho_i(\eta)
&=
\frac{c_i(\eta^{i,0})}{\lambda_i(\eta)}
\qquad\text{when }\lambda_i(\eta)>0.
\end{align*}
When~$\lambda_i(\eta)=0$, the value of~$\rho_i(\eta)$ may be chosen
arbitrarily.

\gm{Section needs to be organized differently. We have a general definition and then we give several examples. It should cotain all the spin systems we consider in the manuscrtipt. }
We also use independent Bernoulli refreshes to perturb the dynamics. For
$\mb s=(s_i)_{i\in\Lambda}\in[0,1]^\Lambda$ and~$\varepsilon\geq0$, define
\begin{equation*}
\Ncal^{\mb s}f
=
\sum_{i\in\Lambda}\bb{E_{i,s_i}-I}f,
\qquad
\Lcal^{\varepsilon,\mb s}
=
\Lcal+\varepsilon\Ncal^{\mb s}.
\end{equation*}
This is again a spin system \gm{more precise, you mean in the sense of~\eqref{eq:spin-generator}?}. If~$\lambda_i^{\varepsilon,\mb s}$ and
$\rho_i^{\varepsilon,\mb s}$ denote its update rate and refresh parameter, then
\begin{equation*}
\lambda_i^{\varepsilon,\mb s}(\eta)
=
\lambda_i(\eta)+\varepsilon, \qquad \rho_i^{\varepsilon,\mb s}(\eta)
=
\frac{\lambda_i(\eta)\rho_i(\eta)+\varepsilon s_i}
{\lambda_i(\eta)+\varepsilon}.
\end{equation*}
\gm{yet another example.}
The constant profile~$\mb 0$ gives independent transitions~$1\to0$. We say
that~$\Lcal$ contains a pure-death component of rate~$\varepsilon>0$ when
\begin{equation}
\label{eq:pure-death-component}
\Lcal=\Lcal^\varepsilon+\varepsilon\Ncal^{\mb 0}
\end{equation}
for another spin-system generator~$\Lcal^\varepsilon$.  
\gm{Not a good notation the epsilon should be on the left side, as here the generator~$\mathcal{L}^\varepsilon$ has nothing to do with~$\varepsilon$.} Equivalently, \gm{the equivalently here is quite dense...}
\begin{equation*}
\gm{c_i(\eta^{i,1})} = \lambda_i(\eta)\bb{1-\rho_i(\eta)}
\geq\varepsilon
\qquad\text{for every }i\in\Lambda\text{ and }\eta\in\Omega,
\end{equation*}
since the left-hand side is the~$1\to0$ transition rate at~$i$.


\gm{explain that we go from generator to semigroup what it means: We have a pathwise process~$\eta(t) \in \Omega$ describe in words what id does. Also write what $P_t$}
Since the update rates are uniformly bounded, the standard Harris construction
gives a Feller process~\cite{Harris1978,Liggett1985}. Its semigroup is
$(P_t)_{t\geq0}$, and for an initial law~$\mu$ we write
$(\mu P_t)(f)=\mu(P_tf)$.


A probability measure~$\pi$ is invariant when
$\pi P_t=\pi$ for every~$t$. We call the process ergodic when
$\mu P_t\Rightarrow\pi$ \gm{what type of convergence? In distribution (i.e. weak)? Needs to be specified.} for every initial law~$\mu$; the limit is then the
unique invariant measure. Uniform exponential ergodicity means that there
is~$\gamma>0$ such that, for every local~$f$ \gm{local is not defined}, some~$C_f<\infty$ satisfies
\begin{equation*}
\sup_{\eta\in\Omega}\abs{P_tf(\eta)-\pi(f)}
\leq C_fe^{-\gamma t}.
\end{equation*}
\gm{This fact is fine but unmotivated....My proposal properly define weak convergence just with local functions and give a remark that local functions is sufficient (usually weak convergence defined via continuous function.} Since~$\Omega$ is compact in the product topology, convergence against local
functions determines weak convergence.

Write~$A\Subset\Lambda$ when~$A$ is finite. For~$A\Subset\Lambda$, set
\begin{equation*}
\chi_A(\eta)=\prod_{i\in A}\eta(i),
\qquad
\chi_\vn=1.
\end{equation*}
The monomials supported on a finite set form a basis for the functions of the
corresponding coordinates \gm{should we be more precise here?}. For a density profile
$\mb p=(p_i)_{i\in\Lambda}\in[0,1]^\Lambda$, let~$\mu_{\mb p}$ be the
Bernoulli product measure with
$\mu_{\mb p}\bb{\eta(i)=1}=p_i$ for every~$i$. For a constant profile,
write~$\mu_p$.

\gm{aslo lets be more precise} For any function~$g$ depending only on the spins in~$N(i)$, write its unique
multilinear expansion as
\begin{equation}
\label{eq:rate-expansion}
g(\eta)
=
\sum_{S\subseteq N(i)}\widehat g(S)\chi_S(\eta).
\end{equation}
We call~$\widehat g(S)$ the M\"obius coefficient of~$g$ at~$S$ \gm{why name moebius coefficient?}.



For~$A\Subset\Lambda$,~$P_t\chi_A(\eta)$ is the expected product of the spins
in~$A$ at time~$t$ when the process starts from~$\eta$.


\gm{Not the best to anticipate Section 4 here...we should find another way to present it.}

Section~\ref{sec:signed-dual} constructs a signed set-valued process
$(A_s,\sigma_s)$ and a potential~$V$ giving a Feynman--Kac representation of
$P_t\chi_A(\eta)$. Writing~$\E_A$ for expectation for the dual started
from~$(A,+)$, we assume throughout that
\begin{equation}
\tag{FK}
\label{eq:fk-integrability}
\E_A\Cb{\exp\bb{\int_0^t V(A_s)\,ds}}<\infty
\qquad
(A\Subset\Lambda,\ t\geq0).
\end{equation}
We refer to~\eqref{eq:fk-integrability} as \emph{Feynman--Kac
integrability}; see~\cite[p.~62]{JansenKurt2014}.
